\documentclass[../../../include/open-logic-section]{subfiles}

\begin{document}

\olfileid[tr]{sth}{replacement}{limofsize}
\olsection{Büyüklük Sınırlaması}

Yerine Koyma için ``içsel'' bir gerekçelendirme sunmaya yönelik belki de en
yaygın girişim şu kavrama dayanır:
\begin{enumerate}
	\item[] \limofsize. Sayısı aşırı fazla olmayan her şey topluluğu bir küme
	oluşturur.
\end{enumerate}
Bu ilke Yerine Koyma'yı hemen doğrular. Sonuçta Yerine Koyma ile oluşturulan
hiçbir küme, oluşturulduğu kümeden daha büyük olamaz. Kesin olarak söylersek:
Yerine Koyma'yı kullanarak $\funimage{\tau}{A} =
\Setabs{\tau(x)}{x \in A}$ kümesini oluşturduğunuzu varsayın; o zaman
$\cardle{\funimage{\tau}{A}}{A}$ olur. Dolayısıyla $A$'nın !!{element}s{}ı bir
kümeyi oluşturamayacak kadar çok değilse görüntüleri de
$\funimage{\tau}{A}$ kümesini oluşturamayacak kadar çok değildir.

Yerine Koyma'yı gerekçelendirmek için \limofsize{} ilkesine başvurmanın açık
güçlüğü, \limofsize{} gibi herhangi bir ilkeyi \emph{henüz} ortaya koymamış
olmamızdır. Üstelik \crefrange{sth:story::chap}{sth:z::chap} içinde
birikimli-yinelemeli küme anlayışına ilişkin öykümüzü anlatırken hiçbir şey
\limofsize{} yönünde bir \emph{ipucu} bile vermedi. Boolos'un bir noktada
``Belki de küme teorisinin `ardında' en az iki düşünce bulunduğu sonucuna
varılabilir'' diye yazmasının nedeni tam da budur
\citeyearpar[p.~19]{Boolos1989}. Bir yanda birikimli-yinelemeli küme anlayışı
çevresindeki fikirlerin $\Z$'yi doğrulaması amaçlanır. Öte yanda
\limofsize{} ilkesinin Yerine Koyma'yı doğrulaması amaçlanır.

Fakat sorun, şimdiye dek \limofsize{} hakkında yalnızca \emph{sessiz}
kalmış olmamız değildir. Daha çok, \limofsize{} ilkesinin (az önceki
biçimiyle) birikimli-yinelemeli küme kavramıyla oldukça kötü uyuşmasıdır.
Sonuçta kümelerin \emph{aşamalar} hâlinde oluşturulması fikrinden hiç söz
etmez.
% [tr source emendation] Upstream has `the issue it is not just`; Turkish
% renders the evident intended construction `the issue is not just`.

Bu ``iki düşüncenin'' (yani birikimli-yinelemeli küme anlayışı ile
\limofsize{} ilkesinin) tarihi göz önüne alındığında bu hiç şaşırtıcı
değildir. Bu ``iki düşünce'', nihayetinde küme kuramsal paradoksları engellemeye
yönelik oldukça farklı iki tasarıya karşılık gelir. Birikimli-yinelemeli küme
kavramı Russell paradoksunu kabaca şöyle engeller: \emph{Bir Russell
kümesinin var olmasını hiçbir zaman beklememeliydik; çünkü o hiçbir aşamada
``oluşturulamazdı''.} Buna karşılık \limofsize{} ilkesinin Russell kümesini
kabaca şöyle dışlaması amaçlanır: \emph{Bir Russell kümesinin var olmasını
hiçbir zaman beklememeliydik; çünkü o aşırı büyük olurdu.}

Böyle söylendiğinde açık konuşalım: paradokslara bir yanıt olarak ele alınan
\limofsize{} ilkesinin \emph{çok} daha fazla gerekçelendirilmeye ihtiyacı
vardır. Örneğin Russell Paradoksu'nun şu sürümünü düşünün: \emph{Hiçbir pug,
kendisini koklamayan pugların tam olarak hepsini koklamaz}
(\olref[sth][story][rus]{sec} kesimine bakın). ``Böyle bir pug neden yok?''
diye sorarsanız, onun aşırı sayıda pugı koklamak zorunda olacağının söylenmesi
iyi bir yanıt değildir. Öyleyse bir Russell kümesinin var olmamasına ilişkin,
var olamayacak kadar ``büyük'' olması neden iyi bir sezgisel açıklama olsun?

Kısacası \limofsize{} ilkesinin ``sezgisel'' güdülenimi konusunda biraz
şaşkınsanız bu bağışlanabilir bir durumdur.

\end{document}
